% ===== 原文 pp.8--9 =====
% tr-01 片 c：§2.1.3 末（对 $\alpha$ 的依赖）、§2.1.4 $q$-形变
% 编号公式 tags: 2.19--2.30

% [原文 p.8]
% 注：原文 p.7 末「When the surface has punctures, there」一句跨页，其整句排在下面图 5 之后。
\begin{center}
\begin{minipage}{\textwidth}
\centering
\includegraphics[width=0.55\textwidth]{fig5.png}\par\vspace{2pt}
\figtitle{5}{该振幅不依赖于如何把一个四孔球面切成两个三孔球面}
\end{minipage}
\end{center}

当曲面带有穿刺时，还存在一个自然的边界项，它使积分等于 $\beta(2-2g-n)$。把这一局部但
拓扑的项 (2.18) 计入后，上面看到的因子 $\alpha^{2-2g-n}$ 就变为 $(e^{\beta}\alpha)^{2-2g-n}$。
% [Tutor Note: p.8, 式 (2.18) 下面关于 $\alpha\to e^{\beta}\alpha$ 的一句，有两处要留意。
%   (1) 按 Gauss--Bonnet 定理，对闭曲面有 $\frac{1}{4\pi}\int d^2x\sqrt{\mathsf{g}}\,R=2-2g$，
%       即 $\int d^2x\sqrt{\mathsf{g}}\,R=4\pi(2-2g)$；原文说它「积分为 $\beta(2-2g)$」，
%       是把 $4\pi$ 并入了 $\beta$ 的定义。
%   (2) 若取欧几里得权重 $e^{-S}$，则加入 $\delta S=\beta\int\sqrt{\mathsf{g}}\,R$ 后应出现
%       $e^{-\beta(2-2g-n)}$，即 $\alpha\to e^{-\beta}\alpha$；原文写 $e^{\beta}\alpha$，
%       相当于把 $\beta$ 取成拓扑项系数的相反数（或把权重记作 $e^{+S}$）。
%       符号约定不影响结论：拓扑项只是把 $\alpha$ 重新标度为 $e^{\pm\beta}\alpha$。]
形如 (2.18) 的项，可以通过改变量子场论的正规化（regularization）方案和/或重整化
（renormalization）方案而轻易产生。因此，从原始拉格朗日量 (2.1) 的 $\beta=0$ 出发时去追问
我们会得到什么 $\alpha$ 值，并没有太大意义。

\subsubsection{$q$-形变}

至此，我们都是从作用量 (2.1) 出发求解二维杨--米尔斯理论的。我们也可以改从格点表述
出发。具体地说，在二维曲面上画一个足够细的网格。给每条边 $e$ 赋一个取值于群 $G$ 的动力学
变量 $U_e$，给每个面 $f$ 赋一个 Boltzmann 权重
\begin{equation}
  \psi_f=\sum_R e^{-A_fc_2(R)}\tr_R\prod_e U_e
  \tag{2.19}
\end{equation}
其中乘积沿面 $f$ 周围的边 $e$ 取。那么由
\begin{equation}
  Z=\int\prod_e dU_e\prod_f\psi_f
  \tag{2.20}
\end{equation}
定义的路径积分就给出配分函数 (2.17)。
% [Tutor Note: p.8, 式 (2.19)--(2.20)：原文没有交代 (2.19)(2.20) 怎样还原出 (2.17)，要点是：
%   (2.20) 中每条边上的 $\int_GdU_e$ 把两侧的表示指标配对，用
%   $\int_G\chi_R(U)\chi_{R'}(U^{-1})dU=\delta_{RR'}$（同 (2.6)）把相邻两个面对应的 $R$ 认同；
%   每个面随即留下自己的 $e^{-A_fc_2(R)}$，面与面之间则按 $R$ 缩并。把网格加细到三角剖分后，
%   顶点/边/面的计数只通过欧拉示性数 $2-2g-n$ 进入，合并即得 (2.17)。
%   另注意 (2.19) 中 $\prod_eU_e$ 的次序：对普通群 $G$，迹的循环性使乘积从哪条边起算无关紧要；
%   一旦把 $G$ 换成量子群 $G_q$，次序就必须固定，这正是本节余下部分的主题。]
这一理论的一个有趣形变，是规定边变量 $U_e$ 取值于量子群（quantum group）$G_q$ 而不是取值
于普通群 $G$。量子群是把群的矩阵元变成非对易的而得到的。例如，量子群 $\SU(2)_q$ 由考虑
$2\times2$ 矩阵 $U_i{}^j=\begin{pmatrix}\alpha&\beta\\\gamma&\delta\end{pmatrix}$ 给出，其关系为
\begin{equation}
  \alpha\beta=q^{1/2}\beta\alpha,\quad
  \alpha\gamma=q^{1/2}\gamma\alpha,\quad
  \beta\delta=q^{1/2}\delta\beta,\quad
  \gamma\delta=q^{1/2}\delta\gamma,\quad
  \beta\gamma=\gamma\beta
  \tag{2.21}
\end{equation}

% [原文 p.9]
以及
\begin{equation}
  \alpha\delta-q^{1/2}\gamma\beta=\delta\alpha-q^{-1/2}\gamma\beta=1,
  \tag{2.22}
\end{equation}
其复共轭由下式给出
\begin{equation}
  \bar{U}_{\bar{\imath}}^{\bar{\jmath}}
  =\begin{pmatrix}\alpha^*&\beta^*\\\gamma^*&\delta^*\end{pmatrix}
  =\begin{pmatrix}\delta&-q^{1/2}\gamma\\-q^{-1/2}\beta&\alpha\end{pmatrix}.
  \tag{2.23}
\end{equation}

由于矩阵元本身不再对易，要为格点路径积分 (2.20) 中的变量定出正确的次序就有点微妙，但这
是可以做到的 [11,12]。让我们用一个简单的例子来看看这种复杂性是怎么出现的。由上面给出的
$\SU_q(2)$ 矩阵元的显式对易关系，容易验证
\begin{equation}
  \delta_{j\bar{\jmath}}U_i{}^j\bar{U}_{\bar{\imath}}{}^{\bar{\jmath}}=\delta_{i\bar{\imath}}.
  \tag{2.24}
\end{equation}
然而，由于矩阵元的不对易性，我们有
\begin{equation}
  \delta^{i\bar{\imath}}U_i{}^j\bar{U}_{\bar{\imath}}{}^{\bar{\jmath}}\neq\delta^{j\bar{\jmath}}.
  \tag{2.25}
\end{equation}
而实际上我们有
\begin{equation}
  D^{i\bar{\imath}}U_i{}^j\bar{U}_{\bar{\imath}}{}^{\bar{\jmath}}=D^{j\bar{\jmath}}
  \tag{2.26}
\end{equation}
其中
\begin{equation}
  D^{i\bar{\imath}}=\begin{pmatrix}q^{-1/2}&0\\0&q^{1/2}\end{pmatrix}.
  \tag{2.27}
\end{equation}
% [Tutor Note: p.9, 式 (2.24)--(2.27)：这一段说明量子群上「取迹」为什么必须改。
%   (2.24) 表明 $U_i{}^j$ 与 $\bar{U}_{\bar{\imath}}{}^{\bar{\jmath}}$ 互为 $q$-逆：
%   用 $\delta_{j\bar{\jmath}}$ 把上下指标配对后缩并得到 $\delta_{i\bar{\imath}}$。
%   但 (2.25) 表明，若改用 $\delta^{i\bar{\imath}}$ 去配对，由于矩阵元非对易，缩并结果不再是 $\delta^{j\bar{\jmath}}$。
%   补救办法是把 $\delta$ 换成对角矩阵 $D$（(2.27)），(2.26) 才恢复「缩并给出 $D$」的形式；
%   $D$ 的迹 $q^{-1/2}+q^{1/2}$ 就是 (2.28) 中的量子维数。
%   易错点：(2.27) 的对角元次序是「左上 $q^{-1/2}$、右下 $q^{1/2}$」，不要写反。]
因此，未形变 $\SU(2)$ 中自然的组合 $\delta^{i\bar{\imath}}\delta_{i\bar{\imath}}$ 在 $\SU_q(2)$
的表示论中被修改为
\begin{equation}
  \delta^{i\bar{\imath}}\delta_{i\bar{\imath}}=2
  \quad\rightsquigarrow\quad
  D^{i\bar{\imath}}\delta_{i\bar{\imath}}=q^{1/2}+q^{-1/2}
  \tag{2.28}
\end{equation}
右端被称为 $\SU_q(2)$ 二维表示的量子维数（quantum dimension）。

归根结底，由于规范群现在是量子群，配分函数的最终表达式 (2.17) 中唯一的改变，是把维数
$\dim R$ 换成量子维数 $\dim_qR$。对于一般的 $\SU(N)$ 情形，量子维数为
\begin{equation}
  \dim_qR=\tr_R\diag(q^{(N-1)/2},q^{(N-3)/2},\dots,q^{(1-N)/2}),
  \tag{2.29}
\end{equation}
因此 $q$-形变理论的配分函数是
\begin{equation}
  \psi_{A,g}(U_i)=\alpha^{2-2g-n}\sum_R e^{-Ac_2(R)}
  \frac{\prod_i\chi_R(U_i)}{(\dim_qR)^{2g-2+n}}.
  \tag{2.30}
\end{equation}
% [Tutor Note: p.9, 式 (2.29)：取 $R$ 为 $\SU(N)$ 的基础表示，$\diag(\cdot)$ 的迹直接给出
%   $\dim_qR=q^{(N-1)/2}+q^{(N-3)/2}+\cdots+q^{(1-N)/2}$，即量子整数 $[N]_q$；
%   取 $N=2$ 得 $q^{1/2}+q^{-1/2}$，与 (2.28) 右端一致——这是检查 (2.29) 是否抄对的最快办法。
%   注意指数从 $(N-1)/2$ 起每次减 $1$，共 $N$ 项，末项为 $(1-N)/2$。
%   另 (2.30) 与 (2.17) 的唯一差别就是 $\dim R\to\dim_qR$，$\alpha$ 的幂次与指数 $2g-2+n$ 不变。]
已知同样的形变也会从弦论的角度出现 [13]，尽管在那里背后的量子群并不直接可见。
